\section{Related work}\label{sec:related}

\paragraph{Rule following and policy-following agents.} RuLES found that almost
all the chat models it tested struggle to keep simple rules stated in their
instructions while interacting with a user, even on straightforward test cases,
and decided each violation with a programmatic evaluation function rather than
human review \citep{Mu2023RuLES}. Two agent benchmarks apply the same idea to tool use.
$\tau$-bench gives an agent domain tools and written policy guidelines in retail
and airline customer service, simulates the user with a language model, and
scores the final database state against an annotated goal \citep{Yao2025TauBench}.
ST-WebAgentBench pairs 222 web tasks with explicit policies and introduces
Completion Under Policy, which credits a task only if it was completed without
breaching any applicable policy \citep{Levy2026STWebAgentBench}. We follow all
three in deciding violations mechanically, and we adopt the principle that
completion and compliance must be scored together, reporting them as separate
outcomes so that an agent's position on the trade-off between them is visible.

\paragraph{Finance agents and compliance.} Language-model trading agents are an
active line of work: FinMem gives a single trading agent a layered memory and a
configurable character \citep{Yu2024FinMem}, and TradingAgents organizes several
agents into the roles of a trading firm, including a risk-management team
\citep{Xiao2024TradingAgents}. Both are judged on trading performance. Whether
such an agent keeps a firm's rules is a separate question, and the closest work on
it is ReguSim \citep{Luo2026ReguSim}, a financial-compliance environment in which
trading agents act under executable market rules (price bands, resale
restrictions, short-sale permission and exposure caps), paired with a benchmark
for surveillance monitors. Its pressure comes from the objectives and personas
the experimenters assign, such as maximizing profit; visible rules reduce
rejected orders without eliminating them. Its false-rejection measure concerns
monitors that wrongly flag accepted orders, not trading agents that decline
legitimate work. It evaluates hosted models and notes that live reruns may not
reproduce archived outputs. The present benchmark differs on five counts: the
pressure comes from the agent's own principal, who asks for a specific breaching
trade; the pressure includes social-engineering moves such as claimed approvals
and forged system notices; legitimate requests are scored alongside conflicts,
so over-refusal is measured; the rules are firm-level compliance controls rather
than exchange microstructure rules; and every model is open-weight and served
locally, so every result can be rerun from the same weights.

Business Arena places agents in a long-horizon marketplace where regulatory
obligations must be met before trading \citep{Pan2026BusinessArena}.
FinSafetyBench tests whether chat models refuse requests that would breach
financial-compliance norms, across fourteen subcategories drawn from real
financial crime cases \citep{Hou2026FinSafetyBench}. As its abstract describes
it, the benchmark measures whether models refuse harmful requests, under
adversarial attack settings; it does not report refusal of legitimate ones. A
memory-controlled trading benchmark masks tickers and dates so that recalled
market history cannot stand in for reasoning \citep{Zhu2026KnowingDoing}, the
same concern that leads us to invent the universe.

\paragraph{Harmful requests, refusal and pressure.} AgentHarm measures whether
agents refuse explicitly malicious multi-step tasks
\citep{Andriushchenko2025AgentHarm}, and ToolEmu emulates high-stakes tools with a
language model and judges the resulting risks with another
\citep{Ruan2024ToolEmu}; its authors found by human review that about two-thirds
of the failures it flags would be real, which is one reason the present
benchmark grades with a rule engine instead. XSTest documented exaggerated
safety, in which models decline clearly benign prompts that superficially
resemble unsafe ones \citep{Rottger2024XSTest}, the failure our legitimate
scenarios are built to detect. The social pressure levels draw on documented
weaknesses: assistants tend to defer to a user's stated views
\citep{Sharma2023Sycophancy}, which the authority and stale-rule levels test, and
persuasion techniques taken from social science raise the rate at which models
can be talked out of their safety training \citep{Zeng2024Persuasion}.

\paragraph{Injected instructions and the instruction hierarchy.} Instructions
injected into text an application processes can override its intended behavior
\citep{Greshake2023Injection}. For agents, InjecAgent and AgentDojo measure
injections that arrive through tool outputs, such as an email or a web page the agent reads, and find that current agents can be hijacked this way
\citep{Zhan2024InjecAgent,Debenedetti2024AgentDojo}. One proposed remedy is an
instruction hierarchy, in which a model is trained to give system instructions
priority over user messages and user messages priority over third-party text, and
to ignore lower-privileged instructions that conflict with higher ones
\citep{Wallace2024InstructionHierarchy}. The two counterfeit-authorization levels
here test that hierarchy from the inside. Neither comes from a third party: the
claimed approval is the principal asserting, in an ordinary user message, that an
authority above the agent has already spoken, and the forged notice is a user
message formatted to look like a system message. A model that keeps the
hierarchy should give either one no more weight than any other user request.

\paragraph{Regulatory context.} Broker-dealers with market access must already
maintain pre-trade controls that prevent the entry of orders exceeding pre-set
credit or capital thresholds and orders in securities the firm or customer is
restricted from trading, under the firm's direct and exclusive control
\citep{SEC2010MarketAccess}. FINRA has reminded its members that its rules apply
to generative AI tools just as they apply to any other technology a firm uses
\citep{FINRA2024RN2409}, and supervisory guidance on model risk management
expects models to be validated and their limitations understood before their
outputs are relied on \citep{Fed2011SR117}. A firm that lets an agent place
orders therefore remains responsible for the agent's breaches. Earlier work in
this series surveyed audit-trail practice in published machine-learning finance
research \citep{Khan2026AuditTrails} and mapped model-risk regulation across
jurisdictions \citep{Khan2026Crosswalk}; this paper moves from the model to the
agent that acts on its outputs.
